\documentclass[10pt,a4paper]{amsart}
\usepackage[margin=19mm]{geometry}
\usepackage{amsmath,amssymb,mathtools}
\usepackage[scr=rsfs,cal=euler]{mathalfa}
\usepackage{xfrac}
\usepackage[unicode,hidelinks,bookmarks=false]{hyperref}
\newcommand{\ie}{i.e.\ }
\newcommand{\cf}{cf.\ }
\newcommand{\resp}{respectively\ }
\newcommand{\Subsection}{\subsection}
\providecommand{\nobreakdash}{\nobreak\hbox{-}}
\setlength{\emergencystretch}{2em}
\multlinegap0pt
\usepackage{amssymb}
\usepackage{mathtools}
\usepackage[shortlabels]{enumitem}
\usepackage{xcolor}
\newtheorem{proposition}{Proposition}
\title{Codimension-Two Spacelike Submanifolds with Umbilical Lightlike Normal Sections and Their Relationship to Lightlike Hypersurfaces}
\author{Juan S. G\'omez}
\date{Source: arXiv:2506.07934v4 (CC BY 4.0)}
\begin{document}
\maketitle

\noindent\emph{Source and licence.} Codimension-Two Spacelike Submanifolds with Umbilical Lightlike Normal Sections and Their Relationship to Lightlike Hypersurfaces. Version arXiv:2506.07934v4 (CC BY 4.0), distributed by arXiv under the Creative Commons Attribution 4.0 International licence, http://creativecommons.org/licenses/by/4.0/, as recorded in the arXiv licence metadata for this exact version and shown on the abstract page https://arxiv.org/abs/2506.07934v4. Full licence notice: \texttt{licenses/arxiv-2506.07934-CC-BY-4.0.txt}.\par\smallskip

\noindent\textit{Editorial scope.} Counted unit: the article's proposition propo2 (source0.tex lines 436--446). The article's complete original proof of it is delivered with it (source0.tex lines 447--473, one \texttt{proof} environment of 27 lines). No mathematics is written, repaired or re-derived: the statement and the proof are the article's own lines, in the article's own notation, and no retained line is substituted or re-flowed.  Retained uncounted as the source-local context the proof needs: lines 194--196; lines 199--201; lines 296--298; lines 310--312.  Nothing was omitted from any retained span, so the package carries no omission record.  No retained line contains a citation, so the wrapper carries no bibliography and no bibliography entry.  No retained line contains a figure environment, an image-inclusion command or drawing code, so no image is delivered and the package declares no asset dependency.  Editorial formatting only, in addition: an AMS \texttt{amsart} preamble that restates the article's own declarations and macros used by the retained lines, namely the environments \texttt{proposition} on separate counters, together with the standard packages those lines need.  The retained environments print in this document as Proposition 1 in order of appearance, whereas the article prints the same items with the numbering of its own full document.  The complete native source of the article as distributed in the arXiv e-print for this exact version is bundled unchanged as \texttt{sources/179569/source0.tex}.
	\begin{equation} \label{relation}
		g( A_\zeta X, Y) = g( \mathrm{II}(X,Y), \zeta ).
	\end{equation}

	\begin{equation}
		\bigl(\widetilde{R}(X,Y)Z\bigr)^{\top} = R(X,Y)Z + A_{{\rm II}(X,Z)}Y - A_{{\rm II}(Y,Z)}X, \label{GaussEc}
	\end{equation}

	\begin{equation}\label{shear1}
		\mathring{{\rm II}}(X,Y) = {\rm II}(X,Y) - g(X,Y)\mathbf{H},
	\end{equation}

	\begin{equation}\label{shear3}
		g(\mathring{A}_\zeta X, Y) = g(\mathring{\rm II}(X,Y), \zeta),
	\end{equation}

\begin{proposition}\label{propo2}
	Let \(\psi\colon S\to M\) be a codimension-two spacelike submanifold in a spacetime \(M\) that admits an umbilical lightlike normal section. Then,
	\begin{equation}\label{ExtrinsicScalarH}
		\tau_{\mathrm{ext}} = g(\mathbf{H}, \mathbf{H}).
	\end{equation}
	
	Moreover, if \(M\) has constant sectional curvature \(c\), then the scalar curvature of $S$ satisfies
	\begin{equation}\label{Scal}
		\mathrm{Scal} = n(n-1)\bigl(c + g(\mathbf{H}, \mathbf{H})\bigr).
	\end{equation}
\end{proposition}

\begin{proof}
	From~\eqref{shear1}, we have
	\begin{align}
		g(\mathring{\mathrm{II}}, \mathring{\mathrm{II}})
		&= \sum_{i,j=1}^n g\bigl(\mathring{\mathrm{II}}(E_i, E_j),\, \mathring{\mathrm{II}}(E_i, E_j)\bigr) \nonumber \\
		&= \sum_{i,j=1}^n g\bigl(\mathrm{II}(E_i, E_j) - \delta_{ij} \mathbf{H},\, \mathrm{II}(E_i, E_j) - \delta_{ij} \mathbf{H}\bigr) \nonumber \\
		&= \sum_{i,j=1}^n g\bigl(\mathrm{II}(E_i, E_j),\, \mathrm{II}(E_i, E_j)\bigr) - 2\sum_{i=1}^n g\bigl(\mathrm{II}(E_i, E_i),\, \mathbf{H}\bigr) + n\, g(\mathbf{H}, \mathbf{H}) \nonumber \\
		&= g(\mathrm{II}, \mathrm{II}) - n\, g(\mathbf{H}, \mathbf{H}), \label{RelationSecond}
	\end{align}
	where $g(\mathrm{II}, \mathrm{II}) = \sum_{i,j=1}^n g\bigl(\mathrm{II}(E_i, E_j),\, \mathrm{II}(E_i, E_j)\bigr)$, and similarly for $g(\mathring{\mathrm{II}}, \mathring{\mathrm{II}})$. Both $g(\mathrm{II}, \mathrm{II})$ and $g(\mathring{\mathrm{II}}, \mathring{\mathrm{II}})$ are independent of the choice of local orthonormal frame $\{E_1, \ldots, E_n\}$ on $S$.
	
	Therefore, from the Gauss equation~\eqref{GaussEc} and equations~\eqref{relation} and~\eqref{RelationSecond}, we obtain
	\begin{align*}
		g(\mathring{\mathrm{II}}, \mathring{\mathrm{II}}) - n(n-1) g(\mathbf{H}, \mathbf{H}) 
		&= \sum_{i,j=1}^n g\bigl(\mathrm{II}(E_i, E_j), \mathrm{II}(E_i, E_j)\bigr) - \sum_{i,j=1}^n g\bigl(\mathrm{II}(E_i, E_i), \mathrm{II}(E_j, E_j)\bigr) \\
		&= \sum_{i \neq j} g\bigl(\mathrm{II}(E_i, E_j), \mathrm{II}(E_i, E_j)\bigr) - \sum_{i \neq j} g\bigl(\mathrm{II}(E_i, E_i), \mathrm{II}(E_j, E_j)\bigr) \\
		&= -n(n-1) \tau_{\mathrm{ext}}.
	\end{align*}
	This implies that
	\begin{equation} \label{RelationSecondBis}
		g(\mathring{\mathrm{II}}, \mathring{\mathrm{II}}) = n(n-1) \big( g(\mathbf{H}, \mathbf{H}) - \tau_{\mathrm{ext}} \big).
	\end{equation}
	
	Since $\psi$ admits an umbilical lightlike normal section $\xi$, we have $\mathring{A}_\xi = 0$. It then follows from~\eqref{shear3} that, at each point, the total shear tensor takes values in the lightlike direction spanned by $\xi$. Therefore, $g(\mathring{\mathrm{II}}, \mathring{\mathrm{II}}) = 0$, and from~\eqref{RelationSecondBis} we obtain~\eqref{ExtrinsicScalarH}.
	
	In particular, if $M$ is a spacetime of constant sectional curvature $c$, a straightforward computation yields~\eqref{Scal}.
\end{proof}

\end{document}
